\section{Выпуклые задачи поиска минимума функционала}
\label{sec:convex-functional-minimization}

\subsection{Выпуклые множества и выпуклые функционалы}

Пусть $V$ --- нормированное линейное пространство, $K\subset V$ --- множество
допустимых элементов. Множество $K$ называется \emph{выпуклым}, если вместе с любыми
$u,v\in K$ оно содержит весь отрезок, соединяющий эти элементы:
\[
    \lambda u+(1-\lambda)v\in K,
    \qquad 0\leq\lambda\leq 1.
\]

Функционал $J\colon K\to\mathbb R$ называется \emph{выпуклым}, если
\[
    J\bigl(\lambda u+(1-\lambda)v\bigr)
    \leq \lambda J(u)+(1-\lambda)J(v)
\]
для любых $u,v\in K$ и $0\leq\lambda\leq1$. Если при $u\neq v$ и
$0<\lambda<1$ неравенство строгое, то функционал называется строго выпуклым.

Основная задача имеет вид
\[
    \text{найти }u\in K:\qquad
    J(u)=\inf_{v\in K}J(v).
\]
Для выпуклой задачи любой локальный минимум является глобальным. Если $J$ строго
выпуклый, то точка минимума, если она существует, единственна.

\subsection{Коэрцитивность и существование минимума}

Для задач в бесконечномерных пространствах важны не только выпуклость, но и
компактностные свойства. Функционал $J\colon V\to\mathbb R$ называется
\emph{коэрцитивным}, если
\[
    \|v_k\|\to\infty
    \quad\Longrightarrow\quad
    J(v_k)\to+\infty.
\]
Коэрцитивность не позволяет минимизирующей последовательности ``уходить на
бесконечность''.

Функционал называется \emph{слабо полунепрерывным снизу}, если из
$v_k\rightharpoonup v$ следует
\[
    J(v)\leq \liminf_{k\to\infty}J(v_k).
\]
В частности, непрерывный выпуклый функционал в гильбертовом пространстве слабо
полунепрерывен снизу.

Стандартная теорема существования формулируется следующим образом. Если $V$ ---
рефлексивное банахово пространство, $K\subset V$ непусто, выпукло и замкнуто,
а $J$ коэрцитивен и слабо полунепрерывен снизу, то минимум $J$ на $K$
достигается. При строгой выпуклости $J$ решение единственно.

Идея прямого метода вариационного исчисления состоит в выборе минимизирующей
последовательности $v_k$, выделении из неё слабо сходящейся подпоследовательности
и переходе к пределу с использованием слабой полунепрерывности снизу.

\subsection{Производная Гато и условия оптимальности}

Функционал $J$ имеет в точке $u$ \emph{производную Гато}, если существует линейный
непрерывный функционал $J'(u)\in V^*$ такой, что для любого направления $h\in V$
\[
    \langle J'(u),h\rangle
    =\lim_{t\to0}
      \frac{J(u+th)-J(u)}{t}.
\]

Если $J$ выпуклый и дифференцируемый по Гато, то $u\in K$ является точкой
минимума тогда и только тогда, когда выполняется вариационное неравенство
\[
    \langle J'(u),v-u\rangle\geq0,
    \qquad \forall v\in K.
\]
Если ограничений нет, то $K=V$ и это условие принимает привычный вид
\[
    J'(u)=0.
\]
Для выпуклого функционала это условие не только необходимо, но и достаточно.

Важный случай, возникающий в краевых задачах, имеет вид
\[
    J(v)=\frac12 a(v,v)-\ell(v),
\]
где $a(\cdot,\cdot)$ --- симметричная билинейная форма, а $\ell\in V^*$.
Если форма $V$-эллиптична,
\[
    a(v,v)\geq c\|v\|^2,
\]
то $J$ является строго выпуклым и коэрцитивным. Минимизация на выпуклом
множестве $K$ эквивалентна вариационному неравенству
\[
    a(u,v-u)\geq \ell(v-u),
    \qquad \forall v\in K.
\]

\subsection{Недифференцируемые выпуклые функционалы и субградиент}

Для выпуклого, но недифференцируемого функционала используется понятие
\emph{субградиента}. Функционал $p\in V^*$ называется субградиентом $J$ в точке
$u$, если
\[
    J(v)\geq J(u)+\langle p,v-u\rangle,
    \qquad \forall v\in V.
\]
Множество всех субградиентов обозначается $\partial J(u)$ и называется
субдифференциалом. Для дифференцируемого выпуклого функционала
\[
    \partial J(u)=\{J'(u)\}.
\]
В безусловной задаче условие минимума можно записать как
\[
    0\in\partial J(u).
\]

Часто функционал представляют в виде $J=J_0+j$, где $J_0$ выпуклый и
дифференцируемый, а $j$ выпуклый, но может быть недифференцируемым. Тогда
минимизация на $K$ эквивалентна условию
\[
    \langle J_0'(u),v-u\rangle+j(v)-j(u)\geq0,
    \qquad \forall v\in K.
\]

\subsection{Основные численные подходы}

После конечномерной аппроксимации задача сводится к выпуклой оптимизации.
Используются градиентные методы, метод Ньютона, метод условного градиента,
проекционные методы. Ограничения можно учитывать непосредственно либо заменять
задачу последовательностью задач без ограничений с помощью штрафных и
барьерных функционалов.

Для недифференцируемых выпуклых задач применяются субградиентные и
проксимационные методы. В гильбертовом пространстве проксимальное приближение
строится как минимум строго выпуклого функционала вида
\[
    \frac12\|v-u_k\|^2+\alpha_k J(v).
\]
Квадратичный член стабилизирует задачу и обеспечивает единственность очередного
приближения при стандартных предположениях.


\subsection{Первое условие выпуклости и сильная выпуклость}

Для дифференцируемого выпуклого функционала справедливо опорное неравенство
\[
    J(v)\ge J(u)+\langle J'(u),v-u\rangle,
    \qquad u,v\in K,
\]
если отрезок между $u$ и $v$ содержится в области определения. Именно из него
следует достаточность условия $J'(u)=0$ в безусловной задаче.

Функционал называется $\mu$-сильно выпуклым, если существует $\mu>0$ такое, что
\[
J(v)\ge J(u)+\langle J'(u),v-u\rangle
      +\frac{\mu}{2}\|v-u\|^2.
\]
Сильная выпуклость обеспечивает единственность минимума и даёт количественную
оценку отклонения от него:
\[
    J(v)-J(u_*)\ge \frac{\mu}{2}\|v-u_*\|^2.
\]

\subsection{Прямой метод: схема доказательства существования}

Пусть
\[
    m=\inf_{v\in K}J(v)
\]
и $\{v_n\}\subset K$ --- минимизирующая последовательность: $J(v_n)\to m$.
Коэрцитивность даёт ограниченность $\{v_n\}$. В рефлексивном банаховом
пространстве из ограниченной последовательности можно выделить слабо сходящуюся
подпоследовательность
\[
    v_{n_k}\rightharpoonup u.
\]
Замкнутое выпуклое множество слабо замкнуто, поэтому $u\in K$. По слабой
полунепрерывности снизу
\[
    J(u)\le \liminf_{k\to\infty}J(v_{n_k})=m.
\]
Но по определению инфимума $J(u)\ge m$, следовательно $J(u)=m$.

\subsection{Субдифференциал и нормальный конус}

Для задачи минимизации выпуклого функционала на замкнутом выпуклом множестве
удобно объединить функционал и ограничение. Нормальный конус к $K$ в точке $u$:
\[
N_K(u)=\{p\in V^*: \langle p,v-u\rangle\le0\ \forall v\in K\}.
\]
Тогда условие оптимальности можно записать в компактной форме
\[
    0\in \partial J(u)+N_K(u).
\]
Для гладкого $J$ это эквивалентно вариационному неравенству
\[
    \langle J'(u),v-u\rangle\ge0\qquad \forall v\in K.
\]

\subsection{Квадратичный функционал и связь с вариационной задачей}

Пусть $V$ --- гильбертово пространство,
\[
    J(v)=\frac12 a(v,v)-\ell(v),
\]
где $a$ симметрична, непрерывна и коэрцитивна:
\[
|a(u,v)|\le M\|u\|\|v\|,\qquad
 a(v,v)\ge \alpha\|v\|^2.
\]
Тогда единственный минимум на всём $V$ удовлетворяет
\[
    a(u,w)=\ell(w)\qquad \forall w\in V.
\]
Это одновременно условие стационарности функционала и слабая постановка
соответствующей краевой задачи. Тем самым задачи выпуклой минимизации непосредственно
связаны с вариационными методами и методом конечных элементов.

\subsection{Что важно сказать на экзамене}

Выпуклая задача --- это минимизация выпуклого функционала на выпуклом множестве.
Для существования минимума в бесконечномерном случае существенны коэрцитивность
и слабая полунепрерывность снизу; строгая выпуклость обеспечивает единственность.
Для дифференцируемого функционала основное условие оптимальности имеет вид
\[
    \langle J'(u),v-u\rangle\geq0,
\]
а без ограничений --- $J'(u)=0$. Для недифференцируемого функционала производная
заменяется субдифференциалом.

\newpage
